\section{A ratchet for the quotient}

The finite-start theorem rests on a simple restriction on decreases of $q_n$.

\begin{lemma}[Forced rebound]
If $\Delta q_n=-1$ and $3q_n\le n+1$, then $\Delta q_{n+1}=1$.
\end{lemma}

\begin{proof}
A down-step gives
\[
 q_{n+1}=q_n-1,
 \qquad
 r_{n+1}=2r_n-q_n+n+1.
\]
The numerator governing the next quotient change is
\[
 2r_{n+1}-q_{n+1}=4r_n-3q_n+2n+3.
\]
An up-step at the next transition occurs when this is at least $n+2$. That condition is
\[
 4r_n\ge3q_n-n-1,
\]
which follows from $r_n\ge0$ and $3q_n\le n+1$.
\end{proof}

\begin{theorem}[Ratchet]
Suppose $3q_k\le k+1$ for every $k\in[u,v]$. Then
\[
 q_k\ge q_u-1\qquad(u\le k\le v+1),
\]
and every decrease of $q$ in this interval is undone by the next step.
\end{theorem}

\begin{proof}
Since $3q_k\le k+1$ forces $q_k\le k$, the bounded-quotient lemma applies at every $k\in[u,v]$; in particular $\Delta q_k\ge-1$, so $q$ falls by at most one per step.

We prove by induction on $k\in[u,v+1]$ a statement stronger than the conclusion:
\begin{equation}
 \text{either}\quad q_k\ge q_u,
 \qquad\text{or}\qquad
 q_k=q_u-1\ \text{ and }\ q_{k+1}=q_k+1.
 \label{eq:ratchet-invariant}
\end{equation}
Both alternatives give $q_k\ge q_u-1$, so the theorem follows.

At $k=u$ the first alternative holds. Assume \eqref{eq:ratchet-invariant} at some $k\le v$.

If the second alternative holds at $k$, then $q_{k+1}=q_k+1=q_u$, so the first alternative holds at $k+1$.

Otherwise $q_k\ge q_u$. If $q_{k+1}\ge q_u$, the first alternative holds at $k+1$. If instead $q_{k+1}<q_u$, then $q_{k+1}<q_u\le q_k$, so $\Delta q_k=-1$ and $q_{k+1}=q_k-1$; combined with $q_{k+1}<q_u\le q_k=q_{k+1}+1$ this forces
\[
 q_k=q_u,\qquad q_{k+1}=q_u-1.
\]
Because $k\le v$, the forced-rebound lemma applies at $k$ and yields $\Delta q_{k+1}=1$, that is $q_{k+2}=q_{k+1}+1$. So the second alternative holds at $k+1$, completing the induction.
\end{proof}

\begin{remark}
The second alternative in \eqref{eq:ratchet-invariant} cannot be dropped. Tracking only $q_k\ge q_u-1$ would leave open a flat step at level $q_u-1$ followed by a further descent; it is the forced rebound attached to each \emph{arrival} at $q_u-1$ that excludes this, by ensuring the level is never occupied at two consecutive indices. This induction is machine-checked in Lean as \texttt{Conjecture.ratchet}, together with the forced-rebound lemma \texttt{Conjecture.forced\_rebound} and the drop bound \texttt{Conjecture.quotient\_drop\_le\_one} that it uses; the axiom audit reports no use of \texttt{sorryAx}.
\end{remark}
